\begin{recipe}{Pompoenstamp met rookworst}
\kicker[Hoofdgerecht,Vlees,Nederlands,Doordeweeks]{Hoofdgerecht · vlees · Nederlands · doordeweeks}
\index[register]{Pompoenstamp met rookworst}
\index[register]{Pompoen}
\index[register]{Wortel}
\index[register]{Rookworst}
\index[register]{Stamppot}

\meta{30 minuten}

\lettrine{D}{eze} stamppot maken we met een kant-en-klare zak wortel- en
pompoengroente, gaar gekookt in bouillon en dan gestampt met roomboter.
Net zo makkelijk als hutspot, met een zoetere toon door de pompoen.

\blockrule{Ingrediënten}
\begin{ingredients}
  \ing{1}{rookworst}\index[register]{Rookworst}
  \ing{500 g}{kruimige aardappelen, geschild en in stukken}\index[register]{Aardappelen}
  \ing{600 g}{wortel/pompoen stamppot groente (Lidl)}\index[register]{Pompoen}
  \ing{1}{bouillonblokje (rund, kip of groente)}
  \ing{500 ml}{water}
  \ing{25 g}{roomboter}
  \ingb{scheutje melk}
  \ingb{zout en versgemalen peper}
\end{ingredients}

\blockrule{Bereiding}
\tip{Niet te lang stampen: een stamppot met nog wat grove stukken is
lekkerder dan een gladde puree.}
\begin{steps}
  \item Doe de aardappelen, de wortel/pompoengroente en het bouillonblokje
        in een grote pan. Voeg het water toe en breng aan de kook.
  \item Zet het vuur lager en kook zachtjes ca.\ 20 minuten tot de groenten
        gaar zijn. Verwarm de rookworst ondertussen ca.\ 10 minuten in een
        aparte pan met wat water op laag vuur.
  \item Giet de groenten af.
  \item Stamp grof met de roomboter: er mogen nog wat grove stukken in
        zitten. Voeg een scheutje melk toe voor de gewenste dikte.
  \item Breng op smaak met zout en peper. Serveer met de rookworst erbij.
\end{steps}
\end{recipe}
